\begin{table*}[!t]
\caption{Digital-Mode Signal Classes Included in the Dataset}
\label{tab:signal_modes}
\centering
\footnotesize
\setlength{\tabcolsep}{3pt}
\renewcommand{\arraystretch}{1.12}
\begin{tabular}{clp{0.16\textwidth}p{0.17\textwidth}p{0.38\textwidth}}
\hline
\hline
ID & Dataset Label & Fldigi Modem Name & Signal Family & Main Parameter or Description \\
\hline
0 & BPSK125 & BPSK125 & PSK & Binary phase-shift keying, 125 baud \\
1 & BPSK31 & BPSK31 & PSK & Binary phase-shift keying, 31.25 baud \\
2 & BPSK63 & BPSK63 & PSK & Binary phase-shift keying, 62.5 baud \\
3 & Contestia\_4\_250 & Cont-4/250 & MFSK text mode & 4 tones, 250 Hz nominal bandwidth \\
4 & Contestia\_8\_500 & Cont-8/500 & MFSK text mode & 8 tones, 500 Hz nominal bandwidth \\
5 & CW & CW & On-off keying & Morse-code continuous-wave audio tone \\
6 & DominoEX\_8 & DOMEX8 & Incremental-frequency MFSK & DominoEX 8 mode \\
7 & MFSK16 & MFSK16 & MFSK & Multi-frequency shift keying, 16-tone mode \\
8 & MFSK32 & MFSK32 & MFSK & Multi-frequency shift keying, 32-tone mode \\
9 & Olivia\_4\_250 & OLIVIA-4/250 & MFSK text mode & 4 tones, 250 Hz nominal bandwidth \\
10 & Olivia\_8\_500 & OLIVIA-8/500 & MFSK text mode & 8 tones, 500 Hz nominal bandwidth \\
11 & QPSK125 & QPSK125 & PSK & Quadrature phase-shift keying, 125 baud \\
12 & QPSK31 & QPSK31 & PSK & Quadrature phase-shift keying, 31.25 baud \\
13 & QPSK63 & QPSK63 & PSK & Quadrature phase-shift keying, 62.5 baud \\
14 & RTTY & RTTY & FSK & Two-tone frequency-shift keying radioteletype \\
15 & Thor\_16 & THOR16 & MFSK text mode & THOR 16 mode \\
16 & Throb\_2 & THROB2 & Multi-tone narrowband mode & THROB 2 mode \\
\hline
\hline
\end{tabular}
\end{table*}
